% Source: https://artofproblemsolving.com/community/c3h2958038_triangles_formed_by_diagonals_of_a_polygon
% Stable identity: aops-topic-2958038-q01
We need $3$ vertices with no two adjacent around the cycle. Fix whether vertex $1$ is used.
\begin{asy}
usepackage("amsmath");
size(3cm);
defaultpen(fontsize(10pt));
int n=10;pair[] V;for(int k=0;k<n;++k){V.push(dir(90+36*k));}
for(int k=0;k<n;++k){draw(V[k]--V[(k+1)%n]);dot(V[k]);}
filldraw(V[0]--V[3]--V[6]--cycle,lightgray);
\end{asy}

If vertex $1$ is not used, delete it: we choose $3$ pairwise non-adjacent vertices from a path of $n-1$ vertices, which can be done in $\dbinom{n-3}{3}$ ways. If vertex $1$ is used, its neighbours are excluded, and we choose $2$ non-adjacent vertices from a path of $n-3$ vertices, in $\dbinom{n-4}{2}$ ways. Hence the count is
\[\dbinom{n-3}3+\dbinom{n-4}2=\dfrac{n(n-4)(n-5)}6.\]
For $n=10$ this gives $\dbinom73+\dbinom62=35+15=\boxed{50}$.
